\section*{Chapitre \ref{ch:ml:generative-models-and-synthetic-data} --- Modèles génératifs et données synthétiques}

\begin{solution}{exo:ml:generative-models-and-synthetic-data:1}
\begin{sloppypar}
Sous l'hypothèse nulle, l'exactitude hors échantillon sur 4\,000 fenêtres a pour erreur type $\sqrt{0.25/4\,000} = 0.0079$~: 0.62 se situe à 15
erreurs types au-dessus de 0.5, une différence nette. 0.51 se situe à 1.3 erreur type~: aucun indice de différence, ce qui n'est pas
une preuve d'égalité~; un classifieur plus performant ou davantage d'échantillons pourrait en trouver une.
\end{sloppypar}
\end{solution}

\begin{solution}{exo:ml:generative-models-and-synthetic-data:2}
L'effet de levier et l'asymétrie gains--pertes~: avec des chocs symétriques et une variance qui réagit aux chocs au carré, les rendements négatifs
et positifs font monter la volatilité de la même façon, et la distribution des rendements est symétrique. Un terme GJR ou EGARCH et des chocs dissymétriques
les ajoutent.
\end{solution}

\begin{solution}{exo:ml:generative-models-and-synthetic-data:3}
Un choc sur la variance décroît en $(\alpha + \beta)^k$~; la demi-vie est $\ln 0.5/\ln 0.964 = 19.1$ jours.
\end{solution}

\begin{solution}{exo:ml:generative-models-and-synthetic-data:4}
Cet effet relie le rendement d'un jour à la somme des vingt jours précédents~; à l'intérieur d'un bloc de 20 jours, l'ordre des jours réels est conservé,
si bien que le dernier jour d'une fenêtre et les jours qui le précèdent viennent souvent du même bloc, et la relation y subsiste. Au-delà des
frontières de bloc, elle se rompt. Des blocs de quelques jours la détruisent~; des blocs bien plus longs que vingt jours la conservent presque intacte, et
recopient davantage l'historique.
\end{solution}

\begin{solution}{exo:ml:generative-models-and-synthetic-data:5}
Des fenêtres voisines partagent 31 de leurs 32 jours. Avec un découpage aléatoire, presque chaque fenêtre hors échantillon a un quasi-doublon dans la
moitié d'entraînement portant la même étiquette, si bien que le classifieur reconnaît la période et non la distribution~: 0.83 pour deux moitiés du
même processus. Un découpage par année ne laisse aucun voisin de part et d'autre de la coupure.
\end{solution}

\begin{solution}{exo:ml:generative-models-and-synthetic-data:6}
Le test se déroule dans l'univers du GAN~: une stratégie entraînée et testée sur des trajectoires générées mesure les régularités du générateur,
y compris les artefacts. Ici, un GAN donne 10.7 dans son univers et 0.52 sur des données réelles. Réussir les tests de faits stylisés ne dit rien
sur la prévisibilité. Tester la stratégie sur des données réelles sur lesquelles elle n'a pas été entraînée, et la comparer à un entraînement sur le réel.
\end{solution}

\begin{solution}{exo:ml:generative-models-and-synthetic-data:7}
Après 1\,000 pas~: kurtosis en excès 2.53, exactitude du test à deux échantillons 0.884, ratio de Sharpe TSTS 11.2. Après 8\,000~: 2.49, 0.734 et
9.9 (4\,000 pas ont donné 1.98, 0.749 et 10.7). Un entraînement plus long rend les fenêtres plus difficiles à distinguer en moyenne, mais ne
rétablit pas les queues, que le discriminateur voit rarement sur des fenêtres standardisées, et la prévisibilité inventée subsiste.
Les remèdes sont architecturaux (bruit à queue épaisse, un canal de volatilité comme dans Quant GANs) ou un modèle fondé sur la vraisemblance.
\end{solution}

\begin{solution}{exo:ml:generative-models-and-synthetic-data:8}
$\log p(x) = \log\int p(x\mid z)p(z)\,dz = \log\E_q[p(x\mid z)p(z)/q(z\mid x)]\ge\E_q[\log p(x\mid z)] + \E_q[\log p(z) -
\log q(z\mid x)]$ par l'inégalité de Jensen, et le second terme vaut $-\mathrm{KL}(q\,\|\,p)$. Pour $q = N(m, \mathrm{diag}\,
s^2)$ et $p = N(0, I)$~: $\mathrm{KL} = \tfrac12\sum_j(m_j^2 + s_j^2 - 1 - \log s_j^2)$, le terme du code.
\end{solution}

\begin{solution}{pb:ml:generative-models-and-synthetic-data:1}
\textbf{Partie I.}
\begin{enumerate}
\item Queues épaisses, regroupement de volatilité, un effet de levier, six krachs et un effet de momentum sur 20 jours. Le bootstrap conserve
ce qui tient dans un bloc~; le GARCH-t conserve les queues et le regroupement~; les réseaux peuvent en principe représenter tout cela.
\item $\alpha = 0.079$, $\beta = 0.886$, 3.6 degrés de liberté~; il manque l'effet de levier, le skew et le momentum.
\item \omterm{def:ml:neural-networks-for-noisy-tabular-data:mlp}{Perceptrons multicouches} sur des fenêtres standardisées de 32 jours~: un VAE à code de dimension quatre et un GAN avec un bruit
de dimension 16 et deux couches de 128 unités, 4\,000 pas d'\omterm{def:ml:neural-networks-for-noisy-tabular-data:adam}{Adam} chacun sur un seul cœur.
\item Partir d'un bruit gaussien et appliquer 100 pas de débruitage appris, chacun retirant le bruit prédit et ajoutant un peu
de bruit frais.
\end{enumerate}
\textbf{Partie II.}
\begin{enumerate}[resume]
\item Bootstrap 6, GARCH-t 4, VAE 6, GAN 3, diffusion 4 (réel 6).
\item 0.491, 0.539, 0.619, 0.749, 0.693 (réel contre réel 0.545). Des fenêtres qui se recouvrent fuient à travers un découpage aléatoire~: 0.83 pour
réel contre réel.
\item IC et ratio de Sharpe~: réel 0.078 et 0.86~; bootstrap 0.062 et 0.71~; GARCH-t $-0.019$ et $-0.10$~; VAE 0.040 et 0.76~; GAN
0.052 et 0.52~; diffusion 0.030 et 0.23.
\item Les fenêtres du bootstrap sont proches de fenêtres d'entraînement dans 14.5~\% des cas (copies), celles du GAN dans 11.2~\% (effondrement sur des fenêtres calmes
fenêtres), les autres étant proches des 5~\% de fenêtres de test réelles.
\end{enumerate}
\textbf{Partie III.}
\begin{enumerate}[resume]
\item 10.7, issu de l'autocorrélation linéaire et de profils de volatilité lisses que le GAN a inventés.
\item Le GARCH-t n'a aucune prévisibilité, si bien que la prévision ajustée est une combinaison aléatoire de variables dont l'IC sur données réelles est
aléatoire en signe~: $-0.10$ en moyenne, avec un écart-type de 0.65 sur cinq échantillons.
\item Un test de résistance a besoin de queues et de regroupement~: bootstrap ou GARCH-t, avec les krachs~; l'augmentation a besoin de la
prévisibilité sans artefacts~: aucun des réseaux testés ici n'est assez bon, et les copies du bootstrap apportent peu de nouveauté.
\item Il reproduit des séquences réelles de 20 jours~: quiconque dispose des données réelles peut les reconnaître.
\end{enumerate}
\textbf{Partie IV.}
\begin{enumerate}[resume]
\item \emph{Trajectoires qui ne se sont jamais produites.} Faits stylisés validés et exactitude du test à deux échantillons~: bootstrap 6 et 0.491, GARCH-t 4 et 0.539,
VAE 6 et 0.619, GAN 3 et 0.749, diffusion 4 et 0.693. Ratio de Sharpe TSTR contre 0.86 pour un entraînement sur le réel~: 0.71, $-0.10$, 0.76,
0.52 et 0.23~; l'univers propre du GAN donne 10.7.
\item TSTR proche de l'entraînement sur le réel sur plusieurs échantillons et ajustements, sur une période de test postérieure aux données d'entraînement.
\item Introduire des effets connus dans des données générées, lancer le backtester et vérifier qu'il les retrouve et rien d'autre.
\item Lorsque le mécanisme est en question (impact de marché, dynamique de file d'attente, changement de politique), et lorsque les données sont trop peu nombreuses pour
apprendre une distribution.
\item Fournir le régime en entrée du générateur~; le nouveau risque tient aux régimes rares dans l'entraînement et mal générés
précisément là où les tests de résistance en ont besoin.
\item La version du générateur, la fenêtre des données d'entraînement, ses graines aléatoires, les tests qu'il a réussis, sa finalité, et quels résultats
l'ont utilisé.
\item Les pertes en queue issues des échantillons de diffusion seraient sous-estimées~; il faut faire le test de résistance avec les queues des données réelles ou d'un modèle à
queue épaisse.
\item Sa forme, et seulement autant de sa prévisibilité que l'architecture et les données l'obligent à apprendre.
\end{enumerate}
\end{solution}

\begin{solution}{iq:ml:generative-models-and-synthetic-data:1}
Queues épaisses, absence d'autocorrélation linéaire, regroupement de volatilité, effet de levier, asymétrie gains--pertes, gaussianité
par agrégation, entre autres. Un GARCH à chocs à queue épaisse reproduit les queues, le regroupement et l'agrégation~; les variantes asymétriques
ajoutent l'effet de levier.
\iqlookfor{quatre faits ou plus, et lesquels un GARCH reproduit.}
\end{solution}

\begin{solution}{iq:ml:generative-models-and-synthetic-data:2}
VAE~: un encodeur vers une distribution latente et un décodeur, entraîné sur une borne de vraisemblance~; les échantillons sont flous ou trop gaussiens.
GAN~: un générateur entraîné contre un discriminateur~; \omterm{def:ml:generative-models-and-synthetic-data:gan}{effondrement de modes} (mode collapse) et entraînement instable. Diffusion~: débruitage appris de
données progressivement bruitées~; échantillonnage lent, et les queues dépendent de l'entraînement. Chacun échoue de façon différente sur des données à queues épaisses.
\iqlookfor{le principe d'entraînement et une défaillance pour chacun.}
\end{solution}

\begin{solution}{iq:ml:generative-models-and-synthetic-data:3}
Selon l'objectif~: faits stylisés, un test à deux échantillons avec un découpage par blocs temporels, TSTR contre un entraînement sur le réel sur des données postérieures, et une
vérification de la \omterm{def:ml:large-language-models-in-finance:memo}{mémorisation}~; jamais par les seuls résultats à l'intérieur des données synthétiques.
\iqlookfor{TSTR et un test à deux échantillons par blocs, lié à l'objectif.}
\end{solution}

\begin{solution}{iq:ml:generative-models-and-synthetic-data:4}
Si l'amélioration tient sur des données réelles non utilisées pour ajuster le générateur~; si le générateur a été entraîné sur la
période de test~; si les trajectoires synthétiques portent des artefacts appris par le modèle (TSTS contre TSTR)~; et si le même effet apparaît
avec un bootstrap.
\iqlookfor{évaluation réelle hors échantillon et fuite du générateur.}
\end{solution}

\begin{solution}{iq:ml:generative-models-and-synthetic-data:5}
Le générateur ne couvre qu'une partie de la distribution. Signes dans les rendements~: queues trop fines, trop peu de diversité des
niveaux de volatilité, fenêtres générées proches de quelques fenêtres réelles typiques, un classifieur à deux échantillons qui l'emporte sur les variables de queue.
\iqlookfor{la définition et des diagnostics concrets.}
\end{solution}

\begin{solution}{iq:ml:generative-models-and-synthetic-data:6}
Construire les variables par fenêtre, découper les fenêtres réelles par blocs temporels (avec un écart plus long que la fenêtre) afin qu'aucune fenêtre hors
échantillon ne chevauche une fenêtre d'entraînement, découper les fenêtres générées au hasard, entraîner, et rapporter l'exactitude hors échantillon avec son erreur type.
Un découpage aléatoire place des quasi-doublons de part et d'autre et mesure la reconnaissance des périodes.
\iqlookfor{découpage par blocs avec un écart, et pourquoi le chevauchement gonfle l'exactitude.}
\end{solution}
